\section{Proof automation}\label{automated-sec}

In this project we used proof automation in two ways: automated theorem provers (ATPs) and \emph{Lean} tactics.
ATPs are generally stand-alone tools that implement a (semi-)decision procedure for a given formal language or related set of languages.
For example, \emph{Vampire}~\cite{DBLP:conf/cav/KovacsV13,TheVampireDiary} is an ATP focused primarily on first-order logic using superposition, which we used extensively in this project.  We also made extensive use of \emph{Prover9} and \emph{Mace4}~\cite{prover9-mace4}.

ATPs are complex software that can contain bugs.
Instead of trusting ATP output, we used proof certificates, which many ATPs can produce, to reconstruct proofs in \emph{Lean}.
The details of proof reconstruction depend on the form of the proof certificate produced by the ATP\@.
We expand on this in \Cref{sec:proof-reconstruction}.

Tactics in \emph{Lean}, on the other hand, are meta-programs~\cite{DBLP:journals/pacmpl/EbnerURAM17} that build proofs.
In other words, they essentially take \emph{Lean} code as input and produce \emph{Lean} code as output.
In this manner, they look like another keyword in the language, and are tightly integrated by producing proofs directly.
Under the hood, their implementation can be arbitrarily complex, from syntactic sugar to full decision procedures.
The \texttt{duper} tactic~\cite{DBLP:conf/itp/CluneQBA24}, for example, implements a superposition calculus similar to \emph{Vampire}'s, but for dependent types\,---\,\emph{Lean}'s underlying logical foundation.

In the rest of this section we describe the different proof automation techniques used in this project.
We first discuss the different proof methods used: primarily superposition and equational reasoning; we then discuss the integration in \emph{Lean}, and finally we report some basic empirical results from this project.

\subsection{Proof techniques}

The two main families of ATPs and tactics we used are based on superposition/saturation and equational reasoning.
In this context we also include SMT solvers, which combine specific decision procedures for theories, like congruence closure for equational reasoning, with satisfiability (SAT) solving \cite{deMoura-Bjorner-2009}.
Finally, we also used \texttt{aesop}~\cite{DBLP:conf/cpp/LimpergF23}, which implements a version of tableau search.
This was used mainly to help specific constructions in refutations, and is not specific to proving or disproving magma implications in this sense.
We describe our use of \texttt{aesop} in \Cref{sec:proof-reconstruction} below.

\subsubsection{Saturation}
Most of the ATPs used extensively in this project rely primarily on saturation procedures in the superposition calculus.
For example, this is the case for \emph{Vampire}~\cite{DBLP:conf/cav/KovacsV13}.\footnote{See also~\cite{DBLP:journals/cacm/BentkampBNTVW23} for a gentler exposition.}
The core idea of these provers is that they take a set of assumptions and a conjecture, expressed in (say) first-order logic.
The conjecture is negated and added to the set of assumptions, which are all put into a normal form.
The ATP then tries to refute the negation by applying rules of an underlying calculus, until a proof of false (a contradiction) is derived.
In this case, the conjecture was (classically) true, and the ATP has found a proof by contradiction, often called a ``refutation'' or ``saturation'' proof.

The underlying calculi vary from system to system, but they often have a variant of a resolution clause of the form:
\[\infer{C \lor D}{C \lor L \quad D \lor \neg L} \]
This can be read as $C \lor L$ with $D \lor \neg L$ implies $C \lor D$, where $C, D, L$ are formulas in e.g. first-order logic.
Superposition calculi have a variant of this rule that deals with equality directly, and thus are more efficient at reasoning about equality.

In this project we used \emph{Vampire}~\cite{DBLP:conf/cav/KovacsV13}, \emph{Duper}~\cite{DBLP:conf/itp/CluneQBA24} and \emph{Prover9} and \emph{Mace4}~\cite{prover9-mace4} which are all based on variants of saturation for proving.

\subsubsection{Equational Reasoning} As already discussed in \Cref{canon-sec}, equational reasoning is a type of reasoning that is based\footnote{More precisely, one can formalize this reasoning using Birkhoff's five rules of inference (reflexive, symmetric, transitive, replacement, and substitution); see, e.g., \cite{burris}.} on equational logic and rewriting with congruence~\cite{term-rewriting}.
In general, an equational reasoning procedure takes a series of equations and tries to determine whether another equation can be deduced from it.
A core tool in equational reasoning are \emph{e-graphs}, a data structure used to represent congruence classes of terms.
By themselves, e-graphs provide an efficient means of implementing a decision procedure for congruence closure over ground equations (i.e., equations without variables).
Extensions to this procedure, for example by quantifier instantiation via e-matching \cite{DBLP:conf/cade/MouraB07}, also allow for a semi-decision procedure for congruence closure over non-ground equations.

SMT solvers like \emph{Z3}~\cite{DBLP:conf/tacas/MouraB08} use equational reasoning for deciding the theory of equality with uninterpreted functions~\cite{DBLP:series/txtcs/KroeningS16,DBLP:conf/cade/MouraB07}.
On the other hand, equality saturation~\cite{DBLP:journals/pacmpl/WillseyNWFTP21} uses e-graphs by extending congruence closure to a more controlled search, enabling optimization and conditional rewriting.
One of the main advantages of equational reasoning for implications of magma laws is that we get very explicit proofs: a proof that $l \models l'$ is given by a sequence of rewrites that starts at the left-hand side of $l'$ and arrives at the right-hand side through applications of $l$.

In this project we used \emph{Z3}~\cite{DBLP:conf/tacas/MouraB08}, \emph{Prover9} and \emph{Mace4}~\cite{prover9-mace4}, a custom ATP \emph{MagmaEgg} for magmas based on \emph{egg}~\cite{DBLP:journals/pacmpl/WillseyNWFTP21}, and the \emph{Lean} \texttt{egg} tactic~\cite{DBLP:journals/pacmpl/KoehlerGBGTS24,egg_tactic}, which all work with equational logic. We have also reasoned with manual (custom written) heuristics about simple rewrites.

\subsection{Integration of automation procedures}
\label{sec:proof-reconstruction}

While ATPs are very useful for proving theorems in this project, they do not integrate with \emph{Lean} out of the box.
ATPs may produce unsound proofs, or worse, derive incorrect results.
Thus, by default, theorems in \emph{Lean} cannot be proven by deferring to the result of an ATP\@.
Instead, the results of an ATP can be used to reconstruct a proof of the form required by \emph{Lean}.
Thus, in general, integration of ATPs requires two steps.
First, there is the invocation of the ATPs by translating the problem from \emph{Lean} into the languages and logics they use.
And second, there is the reconstruction of the ATPs' results as a (persistent) \emph{Lean} proof.
These two aspects present different challenges, and require different strategies, depending mostly on the kind of proof strategy the ATP uses.

More generally, we have observed that there are multiple ways of integrating decision procedures within \emph{Lean}, with different levels of integration.

\begin{enumerate}
    \item Using a \emph{Lean} tactic, which calls a decision procedure written in \emph{Lean} (like \texttt{aesop} or \texttt{duper}).
    \item\label{inter} Using a \emph{Lean} tactic, which calls an existing (external) ATP and reconstructs a proof term from the ATP's result (like \texttt{bv\_decide} or \texttt{egg}).
    \item\label{external} Using an external script which calls an existing ATP and generates a source file \texttt{.lean} which captures the result explicitly.
\end{enumerate}

This project primarily used the least integrated approach, Option~\ref{external}, as it was the fastest to implement and imposed no additional technical requirements on other contributors.
The matter of technical requirements caused problems, for example, when integrating the \texttt{egg} tactic (Option~\ref{inter}) as it initially expected certain software on the user's machine.
Such trade-offs between Option~\ref{inter} and Option~\ref{external} are, however, mutual, as the higher upfront cost of integrating a proof tactic in Option~\ref{inter} makes the decision procedure easier to use than with Option~\ref{external}.
Additionally, Option~\ref{inter} can benefit from \emph{Lean}'s meta-programming capabilities when encoding the problem for use with an ATP, and when reconstructing a \emph{Lean} proof from the result.

\subsubsection{Proof Reconstruction}
The relative simplicity of the objects used in this project benefits the implementations of proof reconstruction.
By focusing on the given problem domain, difficult reconstruction issues, like complex dependent types, could be ignored.

For saturation proofs with \emph{Vampire}, we implemented analogs of the \emph{superpose}, \emph{resolve}, and \emph{subsumption} steps in \emph{Lean}.
Proofs can then be reconstructed as sequences of these steps (and additional technicalities) as shown in \Cref{fig:vampire-example}.

\begin{figure}
\centering
\begin{lean}
@[equational_result] theorem Equation999_implies_Equation86
    (G : Type*) [Magma G] (h : Equation999 G) : Equation86 G := by
  by_contra nh
  simp only [not_forall] at nh
  obtain ⟨sK0, sK1, sK2, nh⟩ := nh
  have eq9 (X0 X1 X2 X3 : G) : (X1 ◇ ((X2 ◇ X3) ◇ (X1 ◇ X0))) = X0 :=
    mod_symm (h ..)
  have eq10 : sK0 ≠ (sK1 ◇ (sK2 ◇ (sK1 ◇ sK0))) := mod_symm nh
  have eq11 (X3 X4 X5 : G) : (X4 ◇ (X3 ◇ (X4 ◇ X5))) = X5 := superpose eq9 eq9
  have eq21 : sK0 ≠ sK0 := superpose eq11 eq10
  subsumption eq21 rfl
\end{lean}
\caption{Example of a proof reconstructed from output of \emph{Vampire}. Note how the proof proceeds by contradiction and uses the \texttt{superpose} and \texttt{subsumption} steps implemented in \emph{Lean}.}
\label{fig:vampire-example}
\end{figure}

For equational proofs from external provers, like \emph{MagmaEgg}, we also used a tailored version of reconstruction.
Specifically, the \emph{MagmaEgg} implementation turns \emph{explanations} \cite{nieuwenhuis2005proof} from \emph{egg} into \emph{Lean} proofs by simple applications of the defining properties of equality as shown in \Cref{fig:magma-egg-example}.

\begin{figure}
\centering
\begin{lean}
@[equational_result] theorem Equation3973_implies_Equation4023
    (G : Type _) [Magma G] (h : Equation3973 G) : Equation4023 G := fun x y z =>
  let v0 := M z x
  let v1 := M z v0
  let v2 := M v1 y
  let v3 := M z v1
  have h4 := R v2
  have h5 := R z
  have h6 := h z x v1
  let v7 := M x (M v1 z)
  T (T (T (T (h x y z) (h (M y v0) z v2)) (C (C h5 (C h4 (T (T (h y v0 v2)
  (C (C (T (T (T h6 (h v7 v1 v0)) (C (C (R v1) (T (h v0 v7 z) (C (S h6) h5)))
  (R v0))) (S (h z v1 v0))) (R (M v2 y))) h4)) (S (h y v3 v2))))) h4))
  (S (h (M y v3) z v2))) (S (h v1 y z))
\end{lean}
% SOURCE: https://github.com/teorth/equational_theories/blob/main/equational_theories/Generated/MagmaEgg/small/_005.lean
\caption{Example of a proof reconstructed by \emph{MagmaEgg}. Note the proof only uses reflexivity, symmetry, transitivity, and congruence of equality.}
\label{fig:magma-egg-example}
\end{figure}

In the case of the \texttt{egg} tactic, which also reconstructs proofs from \emph{egg} explanations, the proof could be converted into a more human-readable form by using the \texttt{calcify}\footnote{\url{https://github.com/nomeata/lean-calcify}} tactic, as shown in \Cref{fig:egg-example}

\begin{figure}
  \centering
  \includegraphics[width=\textwidth]{egg-example.png}
  \caption{Example of the \texttt{egg} tactic reconstructing a proof in human-readable form with the help of \texttt{calcify} (invoked by the special syntax \texttt{egg?}).}
  \label{fig:egg-example}
\end{figure}

\subsubsection{Semi-Automated Counterexample Guidance}  Another use of ATPs has been in a semi-automatic fashion, to find counterexamples.
The general strategy was to use ATPs to find counterexamples to implications by building magmas iteratively.
If we want to build a counterexample to $l \models l'$, we want to construct a magma where $l$ holds but $l'$ does not.
In this method, we iteratively strengthen a construction with additional hypotheses, and use the ATP to check whether these hypotheses are not too strong (to imply $l'$) or unsound (to disallow $l$).

% \TODO{this should also be expanded more, at least with references to some of the constructions in other chapters.}

While equational reasoning can also be used in a semi-automatic fashion to prove equations~\cite{DBLP:journals/pacmpl/KoehlerGBGTS24}, the positive implications in the main implication graph of the project were all simple enough that we did not need a semi-automatic approach for them.

\subsection{ATP usage}\label{ATP_usage}

When the project started, contributors had varying degrees of ATP knowledge and skills, with some of us having to start using them from the ground up. With hindsight, several of the project computations could have been approached in better ways, their difficulty diminished due to better ATP expertise. Accordingly, in this section\footnote{Consult the ETP site for a substantially expanded version of these notes.} we provide some facts and hints about ATP usage, primarily with the algebraist working with (unsorted) equational theories in mind, based on our experience and available evidence\footnote{The timings presented here cannot be taken as benchmarks. Different experiments were executed on different machines, with heterogeneous software and OS environments (Ubuntu, Windows 10...), with varying numbers of parallel processes\,---\,both internal and external to the experiment\,---\,and, in many cases, under the Linux \Verb|nice| command.} within the ETP\@. We restrict our attention to \emph{Vampire} and \emph{Prover9}-\emph{Mace4}, which were the main tools used for exploration of implications and anti-implications along the ETP\@.

\subsubsection{Employing several ATPs}
In general, given a batch of problems of interest, it is useful to employ several ATPs when they complement each other on the batch\,---\,i.e., when they can solve different sets of problems. This is often the case with \emph{Vampire} and \emph{Prover9}: e.g., in a 2024 study~\cite{LPAR2024:Prover9_Unleashed_Automated_Configuration}, it is shown that from a batch of around 770 TPTP problems solved with \emph{Vampire} and \emph{Prover9} together (with some restrictions), around $60\%$ are solved by both, $20\%$ are solved only by \emph{Vampire}, and the remaining $20\%$ are solved only by \emph{Prover9}.\footnote{In contrast, in said study the E~prover solves only a single problem not solved by \emph{Vampire}.}

Although older and virtually discontinued, \emph{Prover9} is still very useful with equational logic and algebraic problems. Within the ETP we likewise identified several problems that were easier to solve with \emph{Prover9} than with \emph{Vampire}, and even some solved by \emph{Prover9} but not by \emph{Vampire} (with the configurations we tried). The most salient example is the proof that $\Eq{102744082}$ implies the injectivity of the right multiplication map, used in the Higman-Neumann side project (see \Cref{higman-neumann}), which was originally found by \emph{Prover9} in several hours with parameters chosen to produce a big search space; an optimized choice of \emph{Prover9} options lowers the runtime to 0.2s. Upon contact with the \emph{Vampire} development team, after several attempts they were able to provide a \emph{Vampire} configuration (inspired by the \emph{Prover9} optimization) which produces a proof in 3s.

In addition, \emph{Mace4}'s non-SAT algorithm generally makes it faster at finding models than generic SAT-based finite model builders, such as those implemented in \emph{Vampire}. For ETP problems in particular, a well-configured \emph{Mace4} is faster than \emph{Vampire}'s finite model builder in both the task of finding a model of a given size and that of exhausting a size with no models. For example, for $\Eq{677}$ models, \emph{Mace4} is able to exhaust size 8 in 1.5s, and to find a model of size 9 in 16s, while \emph{Vampire}'s finite model builder is unable to perform any of these tasks in 10 minutes. By contrast, \emph{Vampire}'s \verb|casc_sat| mode can be quicker at finding \emph{some} model of some size for a given problem.

\subsubsection{Input}\label{input}

Both \emph{Prover9}-\emph{Mace4} and \emph{Vampire} present pitfalls regarding quantifier scope, which can lead to mistakes for unwary algebraists.
Let us illustrate this with \emph{Prover9}, using \verb|*| for the operation~$\op$.  In this ATP, free variables are universally quantified at the outermost level of the whole formula, regardless of its parenthesization.
For example, to formulate the property that right multiplications are injective (\texttt{y*x = z*x -> y=z}) if and only they are surjective (\texttt{exists w w*u=v}), one must include universal quantifiers explicitly on both sides of the binary connective \texttt{<->} (quantifiers have higher precedence than \texttt{<->}),
\begin{equation*}
  \texttt{all x all y all z (y*x = z*x -> y=z) <-> all u all v exists w w*u = v.}
\end{equation*}
Omitting the \texttt{all} quantifiers would quantify all variables outside the equivalence, while writing the left-hand side as \texttt{(all x all y all z y*x = z*x -> y=z)} would amount to $(\forall x,y,z, y\op x=z\op x) \implies y'=z'$, with $y'$ and~$z'$ being free variables, hence universally quantified.
We have in fact committed such mistakes over the course of the ETP, thereby giving rise to false proofs of $\Eq{677}\modelsfin\Eq{255}$, our only open implication.

Frequently, it is useful not to run an ATP alone, but to run it in a larger environment permitting multiple calls with different input files, strategies, etc.\@ so that we can provide semi-automated guidance. Integrating the ATP into a computer algebra system further allows to leverage the latter's mathematical capabilities, such as preparing input files with operations from sophisticated algebraic structures. For example, in the ETP we have integrated \emph{Prover9}-\emph{Mace4} with Python and SAGE~\cite{sagemath}, and (among several other applications) used SAGE to access GAP's small groups library~\cite{GAP4,smallgrp} to search, via \emph{Mace4}, for translation-invariant countermodels on specific groups (see \Cref{translation-sec}).

An important factor to consider when generating an input file for proving a conjecture is, for both \emph{Vampire} and \emph{Prover9}, that owing to the way the saturation algorithm operates, the original order in which formulas are presented is actually important: it is perfectly possible to have a collection of axioms which produces a proof in some order but not in another. \emph{Vampire} includes the option \verb|--normalize| to prevent this effect. For \emph{Prover9}, one can use a larger environment to automate the permutations of the input file and make several tries with time restrictions. In contrast, the order of the formulas does not affect \emph{Mace4}'s response.

Moreover, the inclusion of additional formulas redundant with the premises may significantly speed up the proof search, if those formulas are not quickly derived by the ATP algorithm (e.g., evaluations mapping different variables to the same one are routinely included).

Finally, the use of demodulators can greatly simplify and speed up the proof search. In \emph{Prover9}, the \verb|assign(eq_defs, fold)| command allows to substitute any specified subexpression by a user-defined symbol, simplifying all further expressions generated by the ATP\@. For instance, a quick proof that $\Eq{102744082}$ implies the injectivity of right multiplications is obtained by setting \verb|eq_defs| and adding the formula \verb|s(x) = x*x| (where \verb|*| stands for $\op$) to the premises (together with a good weight limit, see \Cref{weight_strategy}). In \emph{Vampire}, deactivating the \verb|--function_definition_elimination| mode can serve a similar purpose.

\subsubsection{Search of compatible properties} When in search of a model, either theoretically or using a model builder, it is useful to be able to identify additional properties to impose on the model, strong enough to simplify the search, yet weak enough to guarantee compatibility with the original axioms. In particular, when seeking a countermodel to implication $\E\models\E'$, we should avoid properties that, together with $\E$, would imply $\E'$. This search can be greatly aided by an ATP: if we have an educated guess that some property $P$ is compatible with the problem $\E\nmodels\E'$, we can run the ATP on $\E\wedge P$ under various strategies to attempt a proof of $\E'$. If, after allowing ample time, no proof is found, this provides heuristic evidence that a countermodel satisfying $P$ exists. Note that even if no proof actually exists, it may well be that there are infinite countermodels but not finite ones.

In the ETP, we successfully applied this approach to several of the outstanding implications. Most notably, we were able to construct an infinite model of $\Eq{1323}\nmodels\Eq{2744}$ after heuristically verifying compatibility\footnote{Specifically, for each property, we ran \emph{Prover9} for 20 minutes with the default configuration and \emph{Vampire} for 999s in CASC mode.} with a unit element and with closure of the operation over the set of square elements.

\subsubsection{Finite setting}\label{finite setting}

There are situations in which one wants to work in the finite setting, for example, when proving a finite implication. Indeed, model builders typically search for finite models; therefore, when using them, we can usually assume we are operating in the finite setting. In this context, injective (resp. surjective) maps are bijective, bijections are periodic maps, and so on. In addition, any law of the form $x = f(y)\op g(x,y)$ for some maps $f,g$ has left multiplication map $L_{f(y)}$ surjective, hence bijective. Thus by applying the substitution $x\mapsto f(y)\op x$ and simplifying $f(y)$ on the left, the law finitely implies $x = g(f(y)\op x,y)$.

By this and similar approaches we typically find some multiplication or related map (squaring, etc.) to be injective and surjective, properties that we may add to our initial formulas. In hindsight, adding injectivity to a finite model builder usually increases performance, whereas surjectivity (and other existentially quantified formulas) tend to diminish it. Moreover, \emph{Prover9}-\emph{Mace4} performs better with operations and equalities than with non-equational formulas. For this reason it is advantageous to convert an injectivity condition from an implication into a new operation. For example, injectivity of the left multiplication map is captured by $x \backslash (x * y) = y$ (where $*$ stands for $\op$); this substitution typically yields a threefold speedup. In contrast, this conversion appears to slow down \emph{Vampire} slightly.

A case study illustrating this technique is given in \Cref{atp-case-study}.
